\section{The intrinsic fibrations of integral homology}

\begin{claim}{E.1}{Refuted}{The proposed comma-category criterion}
The equivalence in Part E.1 is false. In fact, its stated sufficient
conditions do not imply membership in $\Fib_{\WZ}$, even for a map
between finite posets with contractible fibers and comma categories.
\end{claim}
\begin{proof}
Let $B=[2]$, and let $E$ have four objects $a,b,d,c$ and precisely
the nonidentity arrows
\[
                         a\longrightarrow b\longleftarrow d
                           \longrightarrow c.
\]
Define $p(a)=0$, $p(b)=p(d)=1$, and $p(c)=2$. The strict fibers are
$\{a\}$, the interval $d\to b$, and $\{c\}$. The comma categories
$E/0,E/1,E/2$ are respectively $\{a\}$, the span $a\to b\leftarrow d$,
and $E$. Their nerves are trees. Every map in condition (i) is therefore
an integral homology isomorphism. The systems in condition (ii) are
$\mathcal H_0=\mathbb Z$ and $\mathcal H_q=0$ for $q>0$, so (ii)
holds for every integral homology equivalence by its very definition.

Let $v:[1]\to[2]$ select the arrow $0\to2$, and let
$u:[0]\to[1]$ select its initial object. The functor $u$ is a weak
homotopy equivalence over $B$. Its pullback along $p$ is
\[
                          \{a\}\longrightarrow\{a,c\},
\]
where the target is discrete. This map is not an isomorphism on $H_0$.
The file \texttt{code/EH5\_fence\_counterexample.py} specifies all objects,
arrows, and compositions, checks functoriality and the comma categories,
and computes integral homology using integer Smith normal forms.
Its output gives fiber and comma Betti numbers $(1,0,\ldots)$ and
pullback Betti numbers $(2,0)$. Thus the gap precedes the question
about necessity of (ii): the strict-pullback spectral sequence asserted
in the prompt is not justified by its condition (i).
\end{proof}

\begin{definition}[Hypercentral coefficient modules]
For a left $\mathbb ZG$-module $M$, define invariant submodules
transfinitely by
\[
 \begin{split}
 Z_0M&=0,\\
 Z_{\alpha+1}M&=\{m\in M:(g-1)m\in Z_\alpha M
                                      \text{ for every }g\in G\},\\
 Z_\lambda M&=\bigcup_{\alpha<\lambda}Z_\alpha M
                                      \qquad(\lambda\text{ a limit}).
 \end{split}
\]
We call $M$ hypercentral if this filtration exhausts $M$.
The filtration stabilizes by an ordinal of cardinality at most
$|M|^+$, so no proper-class iteration is involved.
\end{definition}

\begin{claim}{E.2}{Proved}{An exact characterization of the coefficient systems}
Let $B$ be a space of CW type, and let $M$ be a local system of abelian
groups on $B$. The following conditions are equivalent:
\begin{enumerate}
\item For every integral homology equivalence $u:Y'\to Y$ over $B$,
the map $H_*(Y';u^*M)\to H_*(Y;M)$ is an isomorphism.
\item For every component and basepoint $b\in B$, the
$\mathbb Z\pi_1(B,b)$-module $M_b$ is hypercentral.
\end{enumerate}
If $\pi_1(B,b)$ is finitely generated, condition (2) is equivalent to
the assertion that for every $m\in M_b$ there is a finite $r$ with
$I^rm=0$, where $I$ is the augmentation ideal. A uniform bound $r$
on all elements is unnecessary. This answers both specified first
cases $B=S^1$ and $B=S^1\vee S^1$.
\end{claim}
\begin{proof}
The submodules $Z_\alpha M_b$ define local subsystems on each component;
each successive quotient has trivial monodromy. An integral homology
equivalence induces an isomorphism with every constant abelian
coefficient group, by the integral universal coefficient exact
sequence. The long exact homology sequences of a short exact sequence
of local systems show that the property in (1) is closed under
extensions. Singular chains with local coefficients commute with
filtered unions of coefficient modules. Transfinite induction,
using this observation at every limit ordinal, proves (2)$\Rightarrow$(1).
Restriction along $Y\to B$ does not affect the argument: the displayed
filtration still has trivial successive quotients after restriction.

For the converse, put $G=\pi_1(B,b)$ and let $Z_\infty M_b$ be the
stable value of the filtration. If $m_0\notin Z_\infty M_b$, stability
implies that some $g_1\in G$ has
$(g_1-1)m_0\notin Z_\infty M_b$. Recursively choose $g_i$ so that
\[
             m_i=(g_i-1)m_{i-1}\notin Z_\infty M_b.
             \tag{E.2.1}
\]
All $m_i$ are nonzero. Let $T$ be a countable rose, with one loop for
each $g_i$, mapped to $B$ by those loops. Put $V=T\vee S^2$.
The universal-cover cellular chain group of $V$ relative to $T$ is a
free rank-one group-ring module in degree two. For each $i$, choose
a map $h_i:V\to V$ fixing $T$, whose action on this relative module
is multiplication by $g_i-1$. Such a map is obtained by pinching the
2-sphere into two spheres, transporting the first around the indicated
loop, and mapping the second with degree $-1$. Equivalently,
$\pi_2(T\vee S^2)$ is the free group-ring module on the attached sphere,
so the element $g_i-1$ supplies the specified attaching map.

The projections $V\to T$ and their composites with $h_i$ are homotopic
relative to $T$, because $T$ is aspherical. Choose these homotopies and
form the mapping telescope $Y=\operatorname{Tel}(h_1,h_2,\ldots)$,
with its resulting map to $T$ and hence to $B$. The inclusion of the
first copy of $T$ gives a map $u:T\to Y$ over $B$. Relative integral
homology of the telescope is the direct limit of $\mathbb Z$ in
degree two under the zero maps, and is therefore zero. Hence $u$ is
an integral homology equivalence.

With the pulled-back system $M$, however, cellular relative chains
give
\[
 H_j(Y,T;M)=
 \begin{cases}
 \colim\bigl(M_b\xrightarrow{g_1-1}M_b
                 \xrightarrow{g_2-1}M_b\longrightarrow\cdots\bigr),&j=2,\\
 0,&j\ne2.
 \end{cases}                                                    \tag{E.2.2}
\]
This follows also from the ordinary mapping-telescope chain complex:
the map $1-\mathrm{shift}$ on the direct sum of the relative chain
groups is injective, and its cokernel is the displayed direct limit.
The class of $m_0$ in this limit is nonzero by (E.2.1). This
contradicts (1), proving hypercentrality. The construction stays in
the fixed universe and uses a countable CW test, regardless of the
cardinality of $G$ or $M_b$.

Suppose now that $G$ has a finite generating set. If every
$(g-1)m$ lies in $Z_\omega M_b$, it suffices to check the finitely many
generators. Their values lie in a common $Z_rM_b$, and the identities
for a product or inverse show the same for every $g\in G$.
Thus $m\in Z_{r+1}M_b$. Consequently $Z_\omega M_b$ is already stable,
and $Z_rM_b=\{m:I^rm=0\}$. This proves the last assertion.

For an explicit failure of uniform nilpotence over the circle, put
\[
 M=\bigoplus_{r\ge1}\mathbb Z[t,t^{-1}]/(t-1)^r.
                                                               \tag{E.2.3}
\]
Every element belongs to finitely many summands and is killed by
some power of $t-1$, whereas no one power kills all summands.
Pull this system back along the map $S^1\vee S^1\to S^1$ that
collapses the second circle for the second requested first case.
\end{proof}

\begin{claim}{E.3}{Proved}{The homology spectral sequence under the correct hypothesis}
For a functor $p:E\to B$, assume that for every $n\ge0$, every
$\sigma:[n]\to B$, and every vertex $i\in[n]$, the strict-fiber
inclusion
\[
             E_{\sigma(i)}\longrightarrow E\times_B[n]
                                                               \tag{E.3.1}
\]
is an integral homology equivalence. Then its fiber homology groups
form local systems $M_q$ on $BB$. For every simplicial set $K\to NB$
there is a natural strongly convergent first-quadrant spectral sequence
\[
       E^2_{s,q}=H_s(|K|;M_q)\ \Longrightarrow\
                H_{s+q}(|NE\times_{NB}K|;\mathbb Z).
                                                               \tag{E.3.2}
\]
There is also the relative sequence for a map $K'\to K$ over $NB$,
with relative homology on both sides. The local system $M_q$ is the
homology of the homotopy fibers of $Bp$ and is canonically identified
with $b\mapsto H_q(E/b)$.
\end{claim}
\begin{proof}
For a simplex $\sigma:\Delta[n]\to K$, put
$P_\sigma=NE\times_{NB}\Delta[n]$. Its chain complex is a covariant
diagram on the simplex category $\Delta/K$. The hypothesis implies
that every map in this diagram induces a homology isomorphism:
compare a vertex of the source simplex with its image vertex in the
target simplex, and apply two out of three to (E.3.1). Hence its
degree-$q$ homology is a local system on $\Delta/K$ and on $|K|$.

The augmentation from the homotopy colimit of the $P_\sigma$ to
$NE\times_{NB}K$ is a weak equivalence. One direct verification uses
the bisimplicial replacement: over a fixed simplex $x$ of the target,
the category of factorizations of its image through simplices of $K$
has the identity factorization as an initial object. Every augmented
row is consequently a contractible resolution of the indicated
simplex. Taking diagonals proves the assertion. The normalized chains
of this replacement form a first-quadrant double complex. Filtering
by the simplex-category degree gives (E.3.2); the homology of its
horizontal complex is precisely homology with the displayed local
system. Taking mapping cones gives the relative version. There are
only finitely many terms on each total-degree diagonal, which proves
strong convergence without any finite-type assumption.

In particular, strict pullback along a weak equivalence of simplicial
sets over $NB$ is an integral homology equivalence, because weak
equivalences preserve homology with every local system. To identify
the fibers, replace $NB$ by a Kan complex $S$ through a weak equivalence,
and pull its path fibration ending at $b$ back to $NB$. The resulting
space $C_b\to NB$ is weakly contractible, and
$NE\times_{NB}C_b$ models the homotopy fiber of $Bp$. The map from the
vertex $b$ to $C_b$ is a weak equivalence over $NB$, so the preceding
observation identifies the homology of the strict and homotopy fibers.
Pullback along the contractible simplex resolution of $B/b$, or
equivalently the same spectral sequence over $N(B/b)$, identifies
this group with $H_q(E/b)$. The construction is natural under an arrow
$b\to b'$ and gives the stated monodromy identification. Finally compare
$NE\to NB$ with an ordinary Kan-fibration replacement over a Kan
replacement of $NB$. On every vertex fiber the comparison is the
homology isomorphism just established. Comparing (E.3.2) with the
ordinary fibration spectral sequence therefore proves that the
strict-to-derived pullback comparison is a homology equivalence for
every base $K\to NB$. Applying the same comparison to mapping cones
proves the relative assertion, with its specified maps.
\end{proof}

\begin{claim}{E.4}{Proved}{An exact intrinsic description of $\Fib_{\WZ}$}
A functor $p:E\to B$ belongs to $\Fib_{\WZ}$ if and only if
(E.3.1) holds and every local system
$M_q(b)=H_q(E_b;\mathbb Z)$ furnished by E.3 is hypercentral.
Thus the characterization uses all finite-simplex pullbacks and
an explicit transfinite module filtration; it does not quantify
over integral homology equivalences of bases.
\end{claim}
\begin{proof}
If $p$ is intrinsic, pull back the weak equivalence
$[0]\xrightarrow{i}[n]$ over each $\sigma$ to obtain (E.3.1).
We explain why its intrinsic property then also applies to
simplicial homology equivalences over $NB$. For $K\to NB$ the
last-vertex functor $\Delta/K\to B$ has nerve a simplicial set
weakly equivalent to $K$ over $NB$. The comparison is the natural
last-vertex map. Its strict pullback along $NE\to NB$ is a homology
equivalence by E.3. If $u:K'\to K$ is a homology equivalence, the
functor $\Delta/u$ belongs to $\WZ$. The intrinsic hypothesis and
these two comparisons imply that its simplicial pullback along
$NE\to NB$ is a homology equivalence as well.

Suppose that some $M_q$ is not hypercentral. Choose the telescope
$u:T\to Y$ of E.2 for that module. It can be represented over $NB$:
choose a simplicial model over a Kan replacement $S$ of $NB$, replace
its map by a Kan fibration, and pull back along $NB\to S$. Right
properness of the ordinary Kan model gives weak equivalences of
these models to the original spaces, compatibly with $u$.
Consequently the preceding paragraph applies. In the relative
spectral sequence (E.3.2), formula (E.2.2), now for every $M_j$,
puts the entire second page in column $s=2$. There are no possible
nonzero differentials. Therefore
\[
 H_{q+2}(E\times_B^hY,E\times_B^hT;\mathbb Z)
   \cong\colim\bigl(M_q\xrightarrow{g_1-1}M_q
                          \xrightarrow{g_2-1}\cdots\bigr)\ne0,
                                                               \tag{E.4.1}
\]
where the notation denotes the pullback models just specified;
the strict models have the same homology by E.3.
This contradicts preservation of $u$. Thus every $M_q$ is
hypercentral. This one-column test is the step that prevents
spectral-sequence differentials from concealing a failure in an
individual fiber homology group.

Conversely, assume (E.3.1) and hypercentrality. For any functor
$u:Y'\to Y$ over $B$ in $\WZ$, E.2 makes the maps of second pages
in (E.3.2) isomorphisms. Strong convergence and the comparison
theorem give an integral homology equivalence
$E\times_BY'\to E\times_BY$. This is exactly the defining
condition for $\Fib_{\WZ}$.
The argument includes empty fibers and disconnected bases; all
coefficient modules are then interpreted componentwise, and the
zero module is hypercentral.
\end{proof}

\begin{claim}{E.5}{Proved}{The trefoil obstruction R4$'$}
If $BB\simeq S^1$, the basepoint map $*\to B$ has no factorization
as an integral homology equivalence followed by an intrinsic
$W_{\mathbb Z}$-fibration. The same assertion holds for homology
with coefficients in any field. In particular the trefoil
obstruction R4$'$ is valid.
\end{claim}
\begin{proof}
Suppose such a factorization has total space $E$, so that $BE$
is acyclic and connected. The finite-simplex consequence of
intrinsic preservation gives the homology spectral sequence
of E.3, also over a field. Let $k$ denote $\mathbb Z$ or a
field, and put $S=\pi_0F$, where $F$ is the homotopy fiber.
Then $H_0(F;k)=k[S]$ is a permutation module for
$\pi_1B=\langle t\rangle$. The edge terms over the circle give
\[
 H_0(E;k)=k[S]_{t},\qquad
 H_1(E;k)\twoheadrightarrow k[S]^{t}.                    \tag{E.5.1}
\]
The first equality makes the action on $S$ transitive. The
second makes its invariant submodule zero. A finite orbit
has the nonzero invariant sum of its basis vectors, over any
of these coefficient rings. Thus $S$ is an infinite orbit,
and $k[S]\cong k[t,t^{-1}]$ with its regular action.

Let $T$ be the trefoil exterior and $a:T\to S^1$ its
abelianization map. The complement has integral homology
that of a circle, for example by Alexander duality, so $a$
is a homology equivalence. Its group has the Wirtinger
presentation
\[
 \pi_1T=\langle x,y\mid xyxy^{-1}x^{-1}y^{-1}\rangle,
 \qquad a_*(x)=a_*(y)=t.
\]
For clarity the required Alexander module is computed here,
rather than inferred merely from an Alexander polynomial.
The infinite cyclic covering cellular complex through
dimension two has
\[
 d_1=(t-1,\ t-1),\qquad
 d_2=\binom{1-t+t^2}{-(1-t+t^2)}.
\]
Since $t-1$ is not a zero divisor, $\ker d_1$ is generated
by $(1,-1)$. Consequently
\[
 H_1(T;k[t,t^{-1}])
       \cong k[t,t^{-1}]/(t^2-t+1)\ne0.                 \tag{E.5.2}
\]
The two Fox derivatives follow by differentiating the six
letters of the relator; they are also checked in
\texttt{code/EH5\_trefoil.py}. The polynomial has nonzero
leading and constant coefficients in every field, so it is
never a Laurent unit. The standard three-crossing Wirtinger
calculation and its Alexander matrix are independently
displayed in \citet[Example 2.35]{EH5KnotNotes2018}.
Only the group presentation is needed for this $H_1$
calculation; no asphericity theorem for knot complements
is being assumed.

After pulling $E\to S^1$ back along $a$, the term
$E^2_{1,0}$ of its spectral sequence is (E.5.2).
No differential enters or leaves this term in a
first-quadrant homology spectral sequence. Thus the
new total space has nonzero $H_1$, whereas $E$ has zero
$H_1$. This contradicts intrinsic preservation of $a$.
Categories of simplices and E.3 supply the same contradiction
for strict categorical pullbacks.

There is also a shorter integral check from E.4:
the regular module $\mathbb Z[t,t^{-1}]$ has no nonzero
invariant elements, hence its transfinite central filtration
is zero. It cannot be one of the fiber modules of an
intrinsic $W_{\mathbb Z}$-fibration.
\end{proof}

\begin{claim}{E.6}{Proved}{The localization-unit criterion R10(a)}
Let $W$ be an accessible fundamental localizer, and let
$\eta_X:X\to LX$ be its localization unit on spaces.
Then $W$ is proper if and only if every homotopy fiber
of every $\eta_X$ is $W$-aspherical.
\end{claim}
\begin{proof}
For an accessible proper localizer, Cisinski's
\citet[Theorem 6.1.11]{EH5Cisinski2006} identifies $W$
with the localizer generated by its aspherical objects.
The associated localization is therefore nullification.
Its local spaces are closed under dependent sums over
a local base: apply $\Map(A,-)$ to such a sum for each
nullifying object $A$, and use
$\Map(A,Z)\simeq Z$ first on the base and then on
every fiber. Thus the units of the associated modality
have $W$-aspherical homotopy fibers. Equivalently, in
the right proper localized simplicial model, pull a
localization unit back along the path fibration of
its local target. That is a localized fibration,
because it is an ordinary fibration between local
spaces. Right properness identifies its total space,
the unit fiber, with a point in $W$.

Conversely let $P$ be the localization generated by all
$W$-aspherical objects. It exists and is accessible
by \citet[Proposition 6.1.8]{EH5Cisinski2006}.
Suppose every unit fiber is $W$-aspherical, hence
$P$-aspherical. For a $P$-local space $Z$, write the
domain of a unit as the homotopy colimit of its
fiber diagram over its full target space. Then
\[
 \Map(X,Z)
   \simeq\holim_{y\in LX}\Map(\fib_y\eta_X,Z)
   \simeq\holim_{y\in LX}Z
   \simeq\Map(LX,Z).                                    \tag{E.6.1}
\]
The middle equivalence is natural in paths and
higher coherences of $LX$; it is not a product
over the set of points. Thus every unit is a
$P$-equivalence. Naturality of units and two out
of three imply $W\subseteq P$, while $P\subseteq W$
by definition. Hence $W=P$, and Theorem 6.1.11
gives properness. This checks the entire
localization-fiber assertion used in R10(a).
\end{proof}

\begin{claim}{E.7}{Proved}{The Morava locally constant localizer is not proper}
For every prime $p$ and every finite height $n\geq1$, the
locally constant localizer of $\operatorname{Mod}_{K(n)}$
is not proper. The witness
\[
 j:K(\mathbb Z,2)\longrightarrow K(\mathbb Z_p,2)          \tag{E.7.1}
\]
in R10(c) is valid, where $\mathbb Z_p$ is the group
of $p$-adic integers regarded discretely.
\end{claim}
\begin{proof}
We use the checked loop criterion: the localizer consists
of maps bijective on components whose loop maps are
$K(n)$-homology equivalences. We first justify the
local target $K(\mathbb Z_p,1)$ without an invalid
inference from mod-$p$ first homology alone.
In \citet[Lemma 7.5]{EH5Bousfield2001}, the localization of
$K(A,m)$ is the modified Postnikov section
$\widehat P^{n+1}_{\tau}(K(A,m)^{\wedge}_p)$.
For $A=\mathbb Z_p$ and $m=1$, this is $K(\mathbb Z_p,1)$
itself: its derived $p$-completion is itself, as follows from
$\mathbb Z_p=\lim_r\mathbb Z_p/p^r$ and the surjectivity of
the reduction maps, and the modified section changes no
group in degree one when $n\geq1$.
Thus $K(\mathbb Z_p,1)$ is $K(n)$-local.
This also repairs the argument from memory in R10(c).
Theorem 1.4 of \citet{EH5BousfieldKn} does establish that a
$K(n)$-equivalence preserves $H_i(-;\mathbb F_p)$ for $i\leq n$,
but that assertion alone does not prove locality with
$\mathbb Z_p$ coefficients. Lemma 7.5 supplies the exact stronger
input required here.

Here is an explicit check of (E.7.1) on loops.
For a torsion-free abelian group $A$,
$H_*(BA;\mathbb Z)=\bigwedge^*A$: express $A$ as
the filtered union of its finitely generated free
subgroups and commute group homology with that union.
The exterior powers are torsion-free. Thus
\[
 H_*(BA;\mathbb F_p)=
             \bigwedge\nolimits_{\mathbb F_p}^*(A/pA).
\]
Since $\mathbb Z/p\cong\mathbb Z_p/p\mathbb Z_p$,
the loop of $j$ is a mod-$p$ homology equivalence.
It is a $K(n)$-equivalence by
\citet[Lemma 10.6]{EH5Bousfield1997}, with
$R_{K(n)}=\mathbb F_p$. Therefore $j\in W^{\mathrm{lc}}_{K(n)}$.

If $A$ is aspherical for this localizer, then
$A$ is simply connected and $\Omega A$ is
$K(n)$-acyclic. Since $K(\mathbb Z_p,1)$ is
$K(n)$-local,
\[
 \Hom(\pi_2A,\mathbb Z_p)=0.
\]
The inclusion $\mathbb Z\hookrightarrow\mathbb Z_p$
also gives $\Hom(\pi_2A,\mathbb Z)=0$. Hurewicz,
the universal coefficient theorem, and simple
connectivity now show that the pointed mapping
spaces from $A$ to both targets in (E.7.1) are
contractible. Hence both targets are local for
the nullification generated by the
$W^{\mathrm{lc}}_{K(n)}$-aspherical spaces.
If the localizer were proper, E.6 would identify
it with that nullification. Its equivalence
$j$ between two local spaces would have to be
a weak equivalence. It is not, since
$\pi_2j$ is the proper inclusion
$\mathbb Z\hookrightarrow\mathbb Z_p$.
\end{proof}

\paragraph{A fixed catalogue for the factorization question.}
Here $L=L_{HZ}$ and $\eta_T:T\to LT$ denotes its unit.
Work first with a simplicial map $f:K\to B$, where $B$ is Kan,
and put
\[
                 \lambda=\max(\aleph_0,|K|,|B|).
\]
The size of a simplicial set means the cardinality of all its
simplices. Fix a set $\mathcal C_\lambda(B)=\{p_i:V_i\to B\}_{i\in I}$
of representatives of the following explicitly specified maps:
\begin{enumerate}
\item $|V_i|\leq\lambda$, and every finite Kan horn lifting problem
for $p_i$ has a solution;
\item for every vertex $b$ and every $q\geq0$, the monodromy module
$H_q((V_i)_b;\mathbb Z)$ satisfies
\[
 Z_{\lambda^+}H_q((V_i)_b;\mathbb Z)=H_q((V_i)_b;\mathbb Z).
                                                               \tag{E.8.1}
\]
\end{enumerate}
There are at most $2^\lambda$ isomorphism types of such maps.
Condition (E.8.1) is a bounded transfinite calculation: every module
in question has cardinality at most $\lambda$, so its upper central
filtration stabilizes before $\lambda^+$. Thus this definition of
the catalogue uses only horn lifting, ordinary fiber homology,
and the filtration of E.2. It makes no reference to an unknown
factorization or to $\Fib_{\WZ}$.

For each catalogue entry form
\[
 P_i=B\times^h_{LB}LV_i,\qquad
 c_i:V_i\longrightarrow P_i,\qquad
 \mathcal A_i(f)=
 \operatorname{Eq}_{/LB}\bigl((LK,Lf),(LV_i,Lp_i)\bigr).
                                                               \tag{E.8.2}
\]
Here $\operatorname{Eq}_{/LB}$ is the space of homotopy equivalences
over $LB$, including the specified homotopy over that base.
For $\alpha\in\mathcal A_i(f)$ the unit and this specified homotopy
give a canonical map $a_{i,\alpha}:K\to P_i$.

The following description uses only fibers of localization units.
For $z\in LT$ write $U_T(z)=\operatorname{hofib}_z(\eta_T)$.
A point $x\in K$ specifies
\[
 z_x=\alpha\eta_K(x),\qquad
 b_x=f(x),\qquad
 \gamma_x:\eta_B(b_x)\simeq Lp_i(z_x).
\]
The map $p_i$ induces
$U_{V_i}(z_x)\to U_B(Lp_i(z_x))$. Define
\[
 D_{i,\alpha}(x)=
 \operatorname{hofib}_{(b_x,\gamma_x)}
 \left(U_{V_i}(z_x)\longrightarrow U_B(Lp_i(z_x))\right).
                                                               \tag{E.8.3}
\]
These spaces, with their transports and higher coherences, form a
family on the fundamental infinity-groupoid of $K$. In particular,
the homotopy limit below retains the specified coherences; replacing
it by a condition on the individual fibers would change the test.

\begin{claim}{E.8}{Proved}{A bounded localization-unit-fiber test for all maps}
The map $f:K\to B$ factors as an integral homology equivalence
followed by a map preserving all integral homology equivalences
under homotopy base change if and only if
\[
 \text{for some }i\in I\text{ and }\alpha\in\mathcal A_i(f),
       \qquad \holim_{x\in K}D_{i,\alpha}(x)\ne\varnothing .
                                                               \tag{E.8.4}
\]
This is a set-sized test: the catalogue is fixed by $B$ and
$\lambda$, independently of $f$, and its membership conditions are
the two displayed conditions above. All remaining obstructions in
(E.8.4) concern maps between the localization-unit fibers in
(E.8.3).

For an arbitrary functor $f:X\to Y$, apply the test to
$NX\to NY\to\operatorname{Ex}^{\infty}NY$ and to
$\lambda=\max(\aleph_0,|NX|,|NY|)$. Then (E.8.4) is equivalent
to the existence of an actual strict factorization
\[
                X\xrightarrow{u}E\xrightarrow{p}Y,
                  \qquad u\in\WZ,\quad p\in\Fib_{\WZ}.
                                                               \tag{E.8.5}
\]
Thus the answer includes all functors, without a fibration or
Grothendieck-construction hypothesis on the given map.
\end{claim}
\begin{proof}
We first prove the size assertion which makes this criterion
bounded. Suppose $K\to V\xrightarrow{p}B$ is a homotopy
factorization of the required kind. Replace $p$ by an ordinary Kan
fibration and represent the first map strictly over $B$. The
Kan-fibration version of E.4, with the same Serre spectral
sequence and telescope proof, says that every
$H_q(V_b;\mathbb Z)$ is hypercentral. Start with the simplicial
subset $V_0\subset V$ generated by the image of $K$. Inductively
construct $V_n\subset V_{n+1}$ by making these three choices in $V$:
\begin{enumerate}
\item For every horn in $V_n$ over a simplex of $B$, include the
simplices of one filler supplied by $p$.
\item For every integral simplicial cycle in $V_n$ which bounds
in $V$, include the finite support of one chain bounding it.
\item For every vertex $b$ of $B$ and every integral simplicial
cycle in $(V_n)_b$ which bounds in $V_b$, include the finite support
of one chain bounding it inside $V_b$.
\end{enumerate}
Take the simplicial closure of all these choices at each stage.
There are at most $\lambda$ choices: horns have finitely many
nondegenerate simplices, integral chains have finite support and
integer coefficients, there are countably many degrees, and
$|B_0|\leq\lambda$. Consequently $|V_n|\leq\lambda$ for every $n$.

Put $V_\infty=\bigcup_nV_n$. Every finite horn occurs at a finite
stage, so $V_\infty\to B$ is Kan. The second and third clauses
give injections
\[
 H_*(V_\infty;\mathbb Z)\hookrightarrow H_*(V;\mathbb Z),
 \qquad
 H_*((V_\infty)_b;\mathbb Z)\hookrightarrow H_*(V_b;\mathbb Z).
                                                               \tag{E.8.6}
\]
Indeed a cycle is supported at one finite stage, and a boundary
in the indicated ambient complex acquires a chosen bounding
chain at the next stage. The total injection is also onto:
every class of $H_*(V)$ comes from $H_*(K)$, and $K\to V$
already factors through $V_\infty$. Thus $K\to V_\infty$ is an
integral homology equivalence. The second injection is a map of
monodromy modules, since both maps to $B$ are Kan fibrations.
A submodule of a hypercentral module is hypercentral: transfinite
induction gives $Z_\beta N=N\cap Z_\beta M$ for $N\subset M$.
Therefore every fiber homology module of $V_\infty$ is
hypercentral. No surjectivity on fiber homology is needed.
We have produced a member of the catalogue through which the
desired factorization passes.

For completeness, an upper central filtration on a set of at
most $\lambda$ elements cannot have a strict successor increase
at every stage below $\lambda^+$. At its first stationary stage
it remains stationary. Hence hypercentrality is exactly
(E.8.1). Labeling the simplices of a set of size at most
$\lambda$ in each degree gives at most $2^\lambda$ choices of
simplicial operators and maps to $B$. This proves all catalogue
bounds, inside the fixed universe.

We next identify the stated obstruction. In any homotopy
commutative square
\[
 \begin{array}{ccc}
 V_i&\xrightarrow{\eta_{V_i}}&LV_i\\
 \big\downarrow p_i&&\big\downarrow Lp_i\\
 B&\xrightarrow{\eta_B}&LB
 \end{array}
\]
the homotopy fiber of
$c_i:V_i\to B\times^h_{LB}LV_i$ at $(b,z,\gamma)$ is
\[
 \operatorname{hofib}_{(b,\gamma)}
       \bigl(U_{V_i}(z)\to U_B(Lp_i(z))\bigr).
\]
This follows directly by writing a homotopy fiber as a point
together with a path to its prescribed image: both sides specify
a point of $V_i$, a path of its localization to $z$, and a
compatible path of its image in $B$ to $b$.
Consequently $D_{i,\alpha}$ is the fiber family of
\[
                K\times^h_{P_i}V_i\longrightarrow K,
                                                               \tag{E.8.7}
\]
and its homotopy limit is precisely the space of homotopy sections
of (E.8.7).

This is also a concrete obstruction test on the cells of $|K|$.
Choose points in the spaces (E.8.3) over the vertices, then
compatible paths over edges. At an $n$-cell, trivialize the
pulled-back family over its contractible characteristic disk.
For $n\geq2$ the boundary lift specifies a class in
$\pi_{n-1}$ of the corresponding space (E.8.3); it must be zero,
and one then chooses a nullhomotopy. For edges the two endpoint
points must lie in the same component after transport.
The choices must extend each other through all skeleta.
Compatible choices produce a section by the CW weak topology,
and a section supplies them. This formulation includes the
inverse-limit compatibility in (E.8.4); nonemptiness of each
separate finite-stage problem would not suffice.

If the original factorization exists, the size argument replaces
it by $K\xrightarrow{u}V_i\xrightarrow{p_i}B$ with
$H_*(u)$ an isomorphism. The map $\alpha=Lu$ is a homotopy
equivalence over $LB$. Naturality of the unit makes $u$ a lift
of $a_{i,\alpha}$ through $c_i$, hence a section of (E.8.7).
This proves necessity.

Conversely, a section of (E.8.7) is a map $u:K\to V_i$ over $B$
together with
$\eta_{V_i}u\simeq\alpha\eta_K$. It may be represented strictly
over $B$, because $K$ is cofibrant and $p_i$ is Kan.
Localizing this identity gives $Lu\simeq\alpha$ under the
canonical unit equivalences. Thus $u$ is an integral homology
equivalence. The defining condition (E.8.1) and the
Kan-fibration version of E.4 imply the required preservation
property for $p_i$. This proves sufficiency for spaces.

We finally verify the assertion about strict functors. The
functorial Kan replacement $\operatorname{Ex}^{\infty}NY$ has
at most $\lambda$ simplices: each finite subdivision of a
simplex is finite, and there are countably many iterations.
An intrinsic factorization of the functor gives a homotopy
factorization over this Kan replacement by E.3--E.4, so the
necessity already proved applies.

In the other direction, pull the simplicial factorization back
along $NY\to B=\operatorname{Ex}^{\infty}NY$. Since $V_i\to B$
is Kan, ordinary right properness shows that
\[
                 V'=NY\times_BV_i\longrightarrow V_i
\]
is a weak equivalence. We obtain a strict simplicial diagram
$NX\to V'\to NY$, with first map an integral homology
equivalence and second map having the required homotopy
preservation property. To realize the original source $X$
itself, use the ordinary proper Thomason model structure,
as established with its natural unit and counit equivalences
in \cite[Theorem~5.2.12]{EH5Cisinski2006}. Applying
$c\operatorname{Sd}^2\operatorname{Ex}^2$ gives a strict diagram
\[
 QX\longrightarrow C\longrightarrow Y,\qquad
 QX=c\operatorname{Sd}^2\operatorname{Ex}^2NX,\quad
 C=c\operatorname{Sd}^2\operatorname{Ex}^2V'.
\]
The counit $q:QX\to X$ is a weak equivalence, and naturality
makes the composite to $Y$ equal to $fq$. The natural
equivalences in that theorem identify the displayed diagram
with $NX\to V'\to NY$ in homotopy.

Factor $QX\to C$, in the slice over $Y$, as an ordinary
Thomason cofibration $QX\to C_0$ followed by an ordinary trivial
fibration $C_0\to C$. Form
$D=X\amalg_{QX}C_0$ in that slice. Left properness makes
$C_0\to D$ a weak equivalence. By two out of three in homology,
the induced map $X\to D$ is an integral homology equivalence,
and $D\to Y$ has the same derived fibers and monodromy as
$V'\to NY$. Replace $D\to Y$ by an ordinary Thomason
fibration $D'\to Y$ through a weak equivalence $D\to D'$.
Right properness makes every vertex-to-simplex pullback map
for $D'\to Y$ a weak equivalence. Its derived fiber homology
is hypercentral, so E.4 gives $D'\to Y\in\Fib_{\WZ}$.
The strict factorization $X\to D'\to Y$ proves (E.8.5).
This argument uses actual slice diagrams throughout.
\end{proof}

\begin{claim}{E.8.1}{Proved}{The boundary cases for the factorization problem}
For connected CW spaces the following assertions hold.
Every map to a connected space $Y$ admits a factorization as an
integral homology equivalence followed by a homotopy map with the
intrinsic preservation property if and only if $H_1(Y;\mathbb Z)=0$.
In particular the basepoint map $*\to Y$ admits such a factorization
if and only if $H_1Y=0$.
If both $X$ and $Y$ are $HZ$-local, then $f:X\to Y$ admits such a
factorization if and only if its homotopy fibers have hypercentral
homology as $\pi_1Y$-modules.
For any factorable map, $H_1f$ is surjective. The converse fails:
the retraction $S^1\vee S^1\to S^1$ is not factorable.
\end{claim}
\begin{proof}
A hypercentral integral permutation module $\mathbb Z[S]$ has
trivial action on $S$. Indeed restrict to the cyclic subgroup
generated by any $g$. Hypercentrality restricts, and E.2 says that
$(g-1)$ acts locally nilpotently. If $gs\ne s$, the elements
$(g-1)^n[s]$ are nonzero for every $n$: for an infinite orbit this
follows by their distinct extreme basis vectors, and for a finite
orbit by tensoring with $\mathbb C$ and diagonalizing its cyclic
permutation, which has a nontrivial root-of-unity eigenvalue.
Thus $gs=s$ for every $g,s$.

For a connected total space $E$ over connected $Y$, the action on
$\pi_0F$ is transitive, with stabilizer the image of
$\pi_1E\to\pi_1Y$. E.4 therefore makes the fiber connected and
this fundamental-group map surjective. In a factorization
$X\xrightarrow{u}E\to Y$ with $u$ a homology equivalence,
$H_1u$ is an isomorphism, proving surjectivity of $H_1f$.
The basepoint case gives the necessary condition $H_1Y=0$.

Suppose conversely that $H_1Y=0$. Attach 2-cells to $Y$ along a
set of generating loops, obtaining a simply connected space $Y_2$.
The relative group $H_2(Y_2,Y)$ is free on those cells. Since
$H_1Y=0$, the map $H_2Y_2\to H_2(Y_2,Y)$ is onto. By Hurewicz
choose based 2-spheres representing lifts of its basis, and attach
3-cells along them. The resulting $j:Y\to Y^+$ has simply
connected target, and its relative cellular complex is the identity
map between the free groups in degrees three and two. Hence $j$
is an integral homology equivalence.
Every simply connected space is $HZ$-local: the groups
$K(A,n)$ are local by the universal coefficient theorem, the
Postnikov tower of a simply connected space is built from them
by homotopy fibers, and local spaces are closed under homotopy
limits. Thus $j$ is an $HZ$-localization unit.
Its homotopy fibers are acyclic. To verify this directly, use the
Serre spectral sequence over the simply connected base $Y^+$.
If $q>0$ were the least degree with nonzero reduced fiber homology,
then the term $E^2_{0,q}$ could only be hit by a transgression
from the base row. Every base-row class survives because
$H_*Y\to H_*Y^+$ is onto in every degree. Thus $E^2_{0,q}$
would give a nonzero kernel of $H_qY\to H_qY^+$, a contradiction.
The fiber is connected by the homotopy exact sequence.

For any $f:X\to Y$, form the canonical homotopy pullback
\[
                  E=Y\times^h_{Y^+}L_{HZ}X,
                                                               \tag{E.8.8}
\]
using the localized map $L_{HZ}X\to Y^+$. The projection
$E\to L_{HZ}X$ has the acyclic fibers just proved, so is a
homology equivalence by the Serre spectral sequence. The
canonical $X\to E$ is consequently a homology equivalence
by two out of three. The other projection $E\to Y$ is a
pullback of a fibration over the simply connected space $Y^+$.
Its fiber homology has trivial monodromy, and E.4 gives the
intrinsic preservation property. This also handles disconnected
$X$, componentwise; an empty source uses the empty total space.

Now assume that both endpoints of $f$ are $HZ$-local, and let
$X\xrightarrow{u}E\xrightarrow{p}Y$ be a requested factorization.
The restriction map
\[
              \Map_Y(E,X)\longrightarrow\Map_Y(X,X),
\]
where these are derived mapping spaces in the slice over $Y$,
is an equivalence. It is the map of homotopy fibers of the two
restriction maps with targets $X$ and $Y$, both equivalences
because these targets are local and $u$ is a homology equivalence.
Thus the identity of $X$ extends to $r:E\to X$ over $Y$, with
$ru\simeq\operatorname{id}_X$ over $Y$. The map $f$ is a
homotopy retract of $p$ in the slice. Preservation of homology
equivalences under homotopy base change is closed under such
retracts, so E.4 applies to $f$. The converse is its ordinary
fibration replacement.

Finally suppose the projection $X=S^1\vee S^1\to S^1$ had a
factorization, with fiber $F$ and total space $E$. The fiber is
connected by the first paragraph. Over the circle the Serre
spectral sequence has only columns zero and one, so
$H_2E=0$ implies $H_1(F)^{t}=0$. A nonzero hypercentral module
has a nonzero invariant element, by taking the first nonzero
stage of its central filtration. Hence $H_1F=0$. The degree-one
sequence then gives $H_1E\cong\mathbb Z$, contradicting
$H_1E\cong H_1X\cong\mathbb Z^2$.
\end{proof}

\begin{claim}{E.9}{Proved}{The locally constant $KU$ localizer is not proper}
Let $W^{\mathrm{lc}}_{KU}$ be the localizer of locally constant
objects valued in $KU$-local spectra. This accessible fundamental
localizer is not proper. An explicit witness to the unit-fiber
criterion E.6 is the localization unit of
\[
                    X=(Q S^5)_{(p)}
                                                               \tag{E.9.1}
\]
for any odd prime $p$. Its homotopy fiber $F$ is simply connected
and has $\pi_2F\cong\mathbb Q/\mathbb Z_{(p)}$, so $F$ is not
$W^{\mathrm{lc}}_{KU}$-aspherical.
\end{claim}
\begin{proof}
Write $L_1$ for localization at the height-one Johnson--Wilson
spectrum $E(1)$, on spaces or spectra as indicated.
It has the same equivalences as $KU_{(p)}$ because
\[
                KU_{(p)}\simeq
                 \bigvee_{i=0}^{p-2}\Sigma^{2i}E(1).
                                                               \tag{E.9.2}
\]
This is the Adams splitting, in the precise form recorded in
\citet[Introduction, displayed splitting]{EH5BakerRichter2005};
smashing the displayed finite wedge with a cofiber proves the
assertion about equivalences.

We record the exact infinite-loop-space localization theorem used.
For a spectrum $V$, let $Z=\Omega^\infty_0V$ and
$Z'=\Omega^\infty_0L_1V$, where the subscript selects the
basepoint component. For a nilpotent space $T$, its modified
Postnikov section $P^2_{\tau}T$ has
\[
 \pi_1P^2_{\tau}T=\pi_1T,\qquad
 \pi_2P^2_{\tau}T=\pi_2T/\operatorname{tors}(\pi_2T),\qquad
 \pi_iP^2_{\tau}T=0\quad(i>2).
\]
The height-one case of
\citet[Theorem 8.1 and definition in Section 7]{EH5Bousfield2001}
is the natural equivalence
\[
 L_1 Z\simeq
      Z'\times^h_{P^2_{\tau}Z'}P^2_{\tau}(Z_{(p)}).
                                                               \tag{E.9.3}
\]
At height one this is an equivalence in every degree, not just
above a range. This last clause is explicit in the cited theorem.

Take $V=\Sigma^4\mathbb S_{(p)}$. Then
$Z=(QS^4)_{(p)}$ is 3-connected, so $P^2_{\tau}Z=*$.
The paragraph preceding Corollary 4.5 in
\citet{EH5Bousfield1979} identifies
$L_1\mathbb S_{(p)}\simeq(L_{KU}\mathbb S)_{(p)}$;
its Corollary 4.5 then gives, for odd $p$,
\[
 \pi_{-3}L_1\mathbb S_{(p)}=0,
 \qquad
 \pi_{-2}L_1\mathbb S_{(p)}
                  =\mathbb Q/\mathbb Z_{(p)}=:C_{p^\infty}.
                                                               \tag{E.9.4}
\]
Consequently $\pi_1Z'=0$ and $\pi_2Z'=C_{p^\infty}$.
The latter group is torsion, so $P^2_{\tau}Z'=*$ also.
Equation (E.9.3) therefore identifies $L_1Z$ with $Z'$,
and in particular
\[
             \pi_1L_1Z=0,\qquad
             \pi_2L_1Z=C_{p^\infty}.                    \tag{E.9.5}
\]

For completeness this is also the $KU$-localization of $Z$,
not merely its $p$-local variant. The target $L_1Z$ is
$KU$-local since every $KU$-equivalence is a
$KU_{(p)}$-equivalence. The source and target are $p$-local
nilpotent spaces. For a prime $q\ne p$ they are mod-$q$
acyclic: their reduced homology groups are $\mathbb Z_{(p)}$-modules,
so multiplication by $q$ is invertible. Lemma 10.6 of
\citet{EH5Bousfield1997} then makes them $KU/q$-acyclic.
Their map is a rational and a $KU/p$ equivalence because it is a
$KU_{(p)}$-equivalence. The arithmetic detection statement in
\citet[Theorem 10.1]{EH5Bousfield1997} makes it a $KU$-equivalence.
Thus (E.9.5) computes $L_{KU}Z$.

Choose a map $f$ presenting $KU$-homology localization.
On connected spaces, the checked loop criterion identifies
$W^{\mathrm{lc}}_{KU}$ with $\Sigma f$-equivalences:
\citet[Theorem 5.2]{EH5Bousfield1997} identifies the specified
unit after looping as
\[
 \Omega X\longrightarrow\Omega L_{\Sigma f}X
            \simeq L_f\Omega X.
                                                               \tag{E.9.6}
\]
Here $\Omega X\simeq Z$. The map induces an isomorphism on
components and, by (E.9.5), both loop spaces are simply connected.
Hence $X$, $L_{\Sigma f}X$, and the homotopy fiber $F$ of their
unit are simply connected. The homotopy exact sequence, using
$\pi_2\Omega X=0$, now gives
\[
             \pi_2F=\pi_1\Omega F
                      \cong C_{p^\infty}.              \tag{E.9.7}
\]
The space $K(C_{p^\infty},1)$ is $L_1$-local by
\citet[Theorem 7.4]{EH5Bousfield2001}: for a connected
Postnikov $H$-space that theorem preserves the $p$-localized
homotopy group in degree one, removes torsion only in degree two,
and rationalizes degrees above two. It is consequently $KU$-local
as well. The nonzero identity homomorphism
$\pi_1\Omega F\to C_{p^\infty}$ represents a nonconstant component
of $\Map_*(\Omega F,K(C_{p^\infty},1))$.
Thus $\Omega F$ is not $KU$-acyclic. The loop criterion says
exactly that $F$ is not $W^{\mathrm{lc}}_{KU}$-aspherical.
Applying E.6 proves nonproperness.

All spaces in (E.9.1)--(E.9.7) admit small simplicial models.
Passing to their categories of simplices therefore supplies the
same localization-unit obstruction for fundamental localizers on
$\Cat$; the conclusion does not rest on a finite-dimensional or
finite-type approximation.
\end{proof}
